\chapter{Interest-Rate Swaps}\label{ch:m2:interest-rate-swaps}

A company has borrowed 200 million dollars for ten years at the overnight rate
plus 1.5\%, and its board wants to know what the interest will cost next year
and every year after that. Its treasurer calls a bank, and the bank quotes a
fixed rate of 4.05\% for ten years: the company will pay that, the bank will pay
the company the overnight rate, and the loan's floating interest is turned
into a fixed 5.55\%. The bank's trader, who has just taken on the other side
of ten years of interest-rate risk, hedges it within the minute, in the
interbank market or on an electronic platform, with an identical swap with
another dealer or with Treasury futures. Interest-rate derivatives make up
about four fifths of the notional of all over-the-counter derivatives, and the
swap curve is a reference against which bank loans, corporate bonds and
structured products are priced and hedged. This chapter
prices a swap, builds the curve it is priced on, and measures its risk.

\section{The swap and its legs}

\begin{definition}[Interest-rate swap, fixed and floating legs]\label{def:m2:interest-rate-swaps:swap}
An \emph{interest-rate swap}\index{interest-rate swap} is an agreement to
exchange, on a notional amount that is never itself exchanged, fixed-rate
interest payments, the \emph{fixed leg}\index{fixed leg}, for floating-rate
payments set by a benchmark rate, the \emph{floating leg}\index{floating leg},
on a schedule of dates until maturity.
\end{definition}

\begin{definition}[Payer and receiver swaps]\label{def:m2:interest-rate-swaps:payer}
In a \emph{payer swap}\index{payer swap} one pays the fixed rate and receives
floating; in a \emph{receiver swap}\index{receiver swap} one receives fixed and
pays floating. A payer gains when rates rise: it is short duration, like a
borrower who has fixed its rate, or like a short position in a bond.
\end{definition}

\begin{omfigure}
\begin{tikzpicture}[font=\small,
  box/.style={draw,thick,rounded corners=2pt,align=center,minimum height=1cm,minimum width=3cm,inner sep=4pt},
  arr/.style={-{Stealth[length=2.5mm]},thick}]
\node[box,draw=omDef,fill=omDef!8] (co) at (0,0) {company\\(payer)};
\node[box,draw=omThm,fill=omThm!8] (bk) at (6.4,0) {bank\\(receiver)};
\node[box,draw=gray,fill=gray!8] (ln) at (-5.2,0) {lenders};
\draw[arr] (co.20) -- node[above=5pt] {fixed 4.05\% a year} (bk.160);
\draw[arr] (bk.200) -- node[below=5pt] {compounded SOFR} (co.340);
\draw[arr] (co) -- node[above=11pt,align=center] {SOFR + 1.50\%} (ln);
\node[font=\footnotesize,text=omThm,align=center] at (3.2,-1.75) {notional USD~200 million: never exchanged,\\only used to compute the payments};
\end{tikzpicture}

\omcaption{The company's hedge. It keeps paying its lenders SOFR plus 1.50\%;
the swap pays it SOFR and charges 4.05\%. The floating payments cancel and the
company pays 5.55\% fixed. Each year's payments are netted and settled two
business days after the period ends.}\label{fig:m2:interest-rate-swaps:flows}
\end{omfigure}

\begin{example}[One year of the company's swap]\label{ex:m2:interest-rate-swaps:period}
The swap starts on 29 September 2026. On 29 September 2027 the company pays
$200\,000\,000 \times 4.05\% \times 365/360 = \text{USD}~8\,212\,500$ and receives the
overnight rate compounded over those 365 days, whatever it turns out to be; on
its loan it pays the same compounded rate plus 1.5\%. The floating payments
cancel, and its cost is fixed.
\end{example}

\section{Overnight index swaps and conventions}

\begin{definition}[Overnight index swap]\label{def:m2:interest-rate-swaps:ois}
An \emph{overnight index swap}\index{overnight index swap} (OIS) is a swap
whose \omterm{def:m2:interest-rate-swaps:swap}{floating leg} pays an \omterm{def:m2:central-banks-and-the-short-rate:rfr}{overnight benchmark rate} \omterm{def:m2:central-banks-and-the-short-rate:arrears}{compounded in arrears} over
each period (\cref{def:m2:central-banks-and-the-short-rate:arrears}). Since the
end of LIBOR it is the standard swap in the currencies whose interbank rate
ceased; in euros it trades alongside swaps on Euribor, which is still
published (\cref{dat:m2:short-term-interest-rate-futures:contracts}).
\end{definition}

\begin{dated}{2026-09}{Conventions of a dollar SOFR swap}\label{dat:m2:interest-rate-swaps:conventions}
Both legs pay annually, actual/360; the \omterm{def:m2:interest-rate-swaps:swap}{floating leg} pays SOFR \omterm{def:m2:central-banks-and-the-short-rate:arrears}{compounded in
arrears}, with payment two government-securities business days after the end of
each period so that the last fixing is known; the swap starts accruing two
business days after the trade date. Conventions differ by currency, and by
the swap's maturity (short swaps may pay once at maturity).
\end{dated}

\section{Pricing, par rates and risk}

\begin{definition}[Annuity and par swap rate]\label{def:m2:interest-rate-swaps:par}
For fixed-leg payment dates $t_1, \dots, t_n$ with accrual fractions $\delta_i$
and discount factors $P(t_i)$, the \emph{annuity}\index{annuity} is $A =
\sum_i \delta_i P(t_i)$, the value of receiving 1 a year on the \omterm{def:m2:interest-rate-swaps:swap}{fixed leg}'s
schedule. The \emph{par swap rate}\index{par swap rate} is the fixed rate at
which the swap is worth zero at inception.
\end{definition}

\begin{proposition}[Value of an OIS in a single curve]\label{prop:m2:interest-rate-swaps:value}
If overnight rates are discounted on the same curve that they project, the
\omterm{def:m2:interest-rate-swaps:swap}{floating leg} of an OIS from $t_0$ to $t_n$ is worth $P(t_0) - P(t_n)$ per unit
of notional, whatever its schedule. Hence the payer's value is
\[
V \;=\; N\bigl[P(t_0) - P(t_n) - K\,A\bigr],
\qquad
K_{\text{par}} \;=\; \frac{P(t_0) - P(t_n)}{A},
\]
and $\partial V/\partial K = -N A$: the swap's sensitivity to its own rate is
its \omterm{def:m2:interest-rate-swaps:par}{annuity}.
\end{proposition}
\begin{proof}
A period's compounded overnight payment at $t_i$ is replicated by investing
$P(t_{i-1})$ at $t_0$ in a deposit to $t_{i-1}$ and rolling 1 overnight from
there, which yields $1/P(t_{i-1}, t_i)$ at $t_i$, of which 1 is returned and the
rest is the payment; its value is $P(t_{i-1}) - P(t_i)$, and the sum telescopes.
\end{proof}

\begin{method}[Bootstrapping an OIS curve]\label{met:m2:interest-rate-swaps:bootstrap}
Take \omterm{def:m2:interest-rate-swaps:par}{par swap rates} for increasing maturities. Solve for the discount factor at
the first maturity so that its swap is worth zero; then, holding it, for the
second, interpolating the discount factors of any intermediate payment dates
between the pillars (here log-linearly in time); and so on. Check that every
input swap reprices at par.
\end{method}

\begin{example}[A ten-year swap]\label{ex:m2:interest-rate-swaps:ten}
With illustrative par rates of 3.90, 3.85, 3.83, 3.85, 3.92 and 4.05\% at one,
two, three, five, seven and ten years, the ten-year par rate is 4.05\% by
construction, the \omterm{def:m2:interest-rate-swaps:par}{annuity} is 8.228 and the swap's DV01 on USD~100 million is
USD~82\,278 by $N A \times 10^{-4}$, and USD~82\,262 when the curve is rebuilt after
bumping each input by a basis point. All of it falls on the ten-year pillar: a
par swap is hedged exactly by the swap that built its pillar.
\end{example}

\begin{omfigure}
\begin{tikzpicture}
\begin{axis}[width=11cm,height=5cm,
  xlabel={maturity (years)},ylabel={rate (\%)},
  tick label style={font=\small},label style={font=\small},
  xmin=0,xmax=10.5,ymin=3.7,ymax=4.5,grid=major,grid style={gray!25},
  legend columns=3,legend style={font=\footnotesize,at={(0.5,-0.3)},anchor=north,draw=none,fill=none,/tikz/every even column/.append style={column sep=8pt}},
  table/col sep=comma]
\addplot[only marks,mark=*,mark size=2pt,omDef] table[x=years,y=par]{figdata/markets-2/09-interest-rate-swaps/par.csv};
\addplot[omThm,very thick] table[x=years,y=zero]{figdata/markets-2/09-interest-rate-swaps/curve.csv};
\addplot[omProp,thick,const plot] table[x=years,y=fwd]{figdata/markets-2/09-interest-rate-swaps/curve.csv};
\legend{par swap rates (inputs),zero rates,one-year forward rates}
\end{axis}
\end{tikzpicture}

\omcaption{The curve bootstrapped from six illustrative par rates. Zero rates
are smooth; forward rates are flat between pillars and jump at them, an
artefact of interpolating the logarithm of discount factors linearly: every
choice of interpolation is a choice of forward curve, and One Quant Book 6
treats better ones. Data: the chapter's tutorial.}\label{fig:m2:interest-rate-swaps:curve}
\end{omfigure}

\begin{omfigure}
\begin{tikzpicture}
\begin{axis}[width=11cm,height=4.8cm,ybar,bar width=12pt,
  ylabel={DV01 (USD thousand per bp)},
  tick label style={font=\small},label style={font=\small},
  xmin=-0.6,xmax=5.6,ymin=-55,ymax=95,grid=major,grid style={gray!25},
  xtick=data,xticklabels from table={figdata/markets-2/09-interest-rate-swaps/buckets.csv}{tenor},
  legend columns=2,legend style={font=\footnotesize,at={(0.5,-0.22)},anchor=north,draw=none,fill=none,/tikz/every even column/.append style={column sep=8pt}},
  table/col sep=comma]
\addplot[fill=omDef!55,draw=omDef] table[x=i,y=par10]{figdata/markets-2/09-interest-rate-swaps/buckets.csv};
\addplot[fill=omThm!55,draw=omThm] table[x=i,y=fwd5y5y]{figdata/markets-2/09-interest-rate-swaps/buckets.csv};
\legend{10-year par swap (payer),5y5y forward swap (payer)}
\end{axis}
\end{tikzpicture}

\omcaption{Bucketed DV01 of two \omterm{def:m2:interest-rate-swaps:payer}{payer swaps} of USD~100 million: the change in
value when each input par rate rises by a basis point and the curve is rebuilt.
The par ten-year loads only its own pillar. The five-year swap starting in five
years is long ten-year rates and short five-year rates: it is a bet on the
forward, and its hedge is a pair of par swaps. Data: the chapter's
tutorial.}\label{fig:m2:interest-rate-swaps:buckets}
\end{omfigure}

\section{Swap spreads}

\begin{definition}[Swap spread]\label{def:m2:interest-rate-swaps:spread}
The \emph{swap spread}\index{swap spread} at a maturity is the \omterm{def:m2:interest-rate-swaps:par}{par swap rate}
minus the yield of the government bond of the same maturity.
\end{definition}

For decades swap rates were above government yields, since the \omterm{def:m2:interest-rate-swaps:swap}{floating leg}
was an interbank rate with bank credit risk in it. In September 2008, shortly
after Lehman Brothers' default, the thirty-year dollar \omterm{def:m2:interest-rate-swaps:spread}{swap spread} turned
negative, and it has stayed negative since. An overnight-rate swap carries
almost no credit risk in its \omterm{def:m2:interest-rate-swaps:swap}{floating leg}, so a negative spread says something
else: owning a Treasury ties up a balance sheet that must be financed in repo,
while receiving fixed on a swap needs only margin, and long-dated demand for
duration, from pension funds hedging their liabilities, meets dealers whose
balance sheets are costly. The \omterm{def:m2:interest-rate-swaps:spread}{swap spread} is the price of that difference,
and trading it, long Treasuries against paying fixed or the reverse, is one of
the core relative-value trades of the rates market (One Quant Book 9).

\section{Execution facilities and compression}

\begin{definition}[Swap execution facility and trade compression]\label{def:m2:interest-rate-swaps:sef}
A \emph{swap execution facility}\index{swap execution facility} (SEF) is a
regulated trading platform for swaps in the United States; swaps that have
been declared available to trade must be executed on one. \emph{Trade
compression}\index{trade compression} is the tearing up of offsetting swaps
between several counterparties, replacing them by fewer swaps with the same
net risk, so as to reduce gross notional, the number of trades and the
capital and margin they consume.
\end{definition}

The swap market was, until the rules that followed the 2008 crisis, largely a
bilateral market between dealers and their clients, traded by telephone. It now has three tiers:
central clearing for most standard swaps (\cref{ch:m2:swap-clearing-and-the-ccp-basis}),
mandatory execution of the most standard ones on electronic platforms, and
reporting of every trade. Compression runs on top: because swaps are never
closed out by selling them, only offset by new swaps, gross notional grows
with every hedge unless it is periodically torn up.

\begin{omfigure}
\begin{tikzpicture}[font=\small,
  node/.style={draw,thick,circle,minimum size=1.1cm,align=center},
  arr/.style={-{Stealth[length=2.5mm]},thick}]
\node[node,draw=omDef,fill=omDef!8] (a) at (0,1.6) {A};
\node[node,draw=omDef,fill=omDef!8] (b) at (2.2,-0.9) {B};
\node[node,draw=omDef,fill=omDef!8] (c) at (-2.2,-0.9) {C};
\draw[arr,omThm] (a) -- node[right=6pt,align=left] {A pays fixed\\to B} (b);
\draw[arr,omThm] (b) -- node[below=4pt] {B pays fixed to C} (c);
\draw[arr,omThm] (c) -- node[left=6pt,align=right] {C pays fixed\\to A} (a);
\node[align=center,anchor=west] at (4.8,0.4) {gross notional before: 300\\net risk of each dealer: 0\\after compression: nothing};
\end{tikzpicture}

\omcaption{A compression cycle. Three dealers hold three identical swaps in a
circle: each pays fixed on one and receives on another, so none has risk, yet
all three hold gross notional, margin and capital against it. Compression
tears the cycle up. Real cycles involve many participants and swaps that
match only approximately; the service finds the largest set of
tear-ups that leaves each participant's risk within its
tolerances.}\label{fig:m2:interest-rate-swaps:compression}
\end{omfigure}

\begin{dated}{2026-09}{The size of the market}\label{dat:m2:interest-rate-swaps:market}
Notional of over-the-counter derivatives outstanding: USD~846 trillion at
end-June 2025 (the latest semiannual BIS survey), 16\% more than a year earlier;
interest-rate derivatives were 79\% of it, about USD~668 trillion, with euro
contracts larger than dollar contracts since 2022. In the United States, SEF
registration was required from October 2013 and execution on a SEF of swaps
declared available to trade from 15 February 2014.
\end{dated}

\section{Tutorial: a curve and a swap}\label{tut:m2:interest-rate-swaps:curve}

\begin{tutorial}
\textbf{Goal.} Bootstrap an OIS curve from six par rates, reprice the inputs,
and compute the bucketed DV01 of a par swap and a forward-starting swap.
\textbf{End state:}
\cref{fig:m2:interest-rate-swaps:curve,fig:m2:interest-rate-swaps:buckets} and the
numbers of \cref{ex:m2:interest-rate-swaps:ten}.
\begin{tutsteps}
\item \textbf{\omterm{def:m2:interest-rate-swaps:par}{Annuity}, par rate, value} (\cref{prop:m2:interest-rate-swaps:value}).
\omcode{code/firm/curve/firm_curve.py}{50}{63}{Annuity, par rate and value of an OIS in a single curve.}\label{lst:m2:interest-rate-swaps:value}
\item \textbf{Bootstrap}: one pillar at a time, by bisection on the logarithm of
the discount factor.
\omcode{code/firm/curve/firm_curve.py}{66}{84}{Bootstrapping discount factors so that each input swap is at par.}\label{lst:m2:interest-rate-swaps:bootstrap}
Every input swap should reprice to zero within a thousandth of a dollar.
\item \textbf{Bucketed DV01}: bump each input, rebuild, reprice.
\omcode{code/firm/curve/firm_curve.py}{87}{95}{Bucketed DV01 by bump and rebuild.}\label{lst:m2:interest-rate-swaps:buckets}
\item \textbf{Run} \texttt{swaps\_demo.ten\_year()} and
\texttt{swaps\_demo.forward\_5y5y()}; \texttt{fig\_swaps.py} writes the charts.
\end{tutsteps}
\textbf{What to change next.} Replace log-linear interpolation of discount
factors by linear interpolation of zero rates and watch the forward curve
change while every input still reprices; then price a swap whose fixed rate is
1\% above par and read its buckets.
\end{tutorial}

\section{Build: the OIS curve and swap pricer}\label{bld:m2:interest-rate-swaps:curve}

\begin{build}
\bfield{Purpose} Every rates position of the miniature firm, and every
discounting of a future cash flow, reads this curve. It is the first curve of
the pricing library (One Quant Book 5 and 6 extend it to several curves, to
turns and to meeting dates, \cref{bld:m2:short-term-interest-rate-futures:meetings}).
\bfield{Interface} \texttt{Curve(spot, times, dfs)} with \texttt{df(date)},
\texttt{zero(date)}; \texttt{schedule(start, years)}; \texttt{annuity};
\texttt{par\_rate}; \texttt{swap\_pv(curve, dates, fixed, notional, payer)};
\texttt{bootstrap(spot, tenors, rates)}; \texttt{bucket\_dv01(\ldots)}.
\bfield{Rules} Annual periods, actual/360 accruals, no business-day adjustment
in this version; log-linear interpolation of discount factors, flat forward
beyond the last pillar; one curve for projection and discounting.
\bfield{Acceptance tests} \texttt{code/firm/curve/tests/}: every input reprices
at par; a flat input curve gives flat par rates at intermediate maturities; a
payer gains when rates rise; a par swap loads only its own pillar; a 5y5y
forward is short the five-year pillar and long the ten-year; leap days.
\bfield{Stretch} Business-day calendars and payment delays (the conventions of
\cref{dat:m2:interest-rate-swaps:conventions}); futures and meeting dates at the
front; a separate discount curve (\cref{ch:m2:swap-clearing-and-the-ccp-basis}).
\end{build}

\begin{omsources}
\item Bank for International Settlements, \emph{OTC derivatives statistics at
end-June 2025}, December 2025.
\item S.\ Klingler and S.\ Sundaresan, ``An explanation of negative swap spreads:
demand for duration from underfunded pension plans'', \emph{Journal of
Finance} 74, 2019.
\item Clarus Financial Technology, ``SOFR swap nuances''; ARRC, \emph{An
Updated User's Guide to SOFR}, 2021.
\item CFTC, final rule on core principles and other requirements for swap
execution facilities, 2013.
\end{omsources}

\section{Exercises}

\begin{exercise}[$\star$]\label{exo:m2:interest-rate-swaps:1}
A \omterm{def:m2:interest-rate-swaps:payer}{payer swap} of USD~100 million at 4.05\% has an annual period of 365 days. What
fixed payment is due at its end? What does the payer receive?
\end{exercise}

\begin{exercise}[$\star$]\label{exo:m2:interest-rate-swaps:2}
A pension fund receives fixed on a thirty-year swap. Is it long or short
duration? Does it gain or lose if rates rise?
\end{exercise}

\begin{exercise}[$\star$]\label{exo:m2:interest-rate-swaps:3}
The ten-year swap rate is 4.05\% and the ten-year Treasury yields 4.20\%. Give
the \omterm{def:m2:interest-rate-swaps:spread}{swap spread} and say what it would have meant before 2008.
\end{exercise}

\begin{exercise}[$\star\star$]\label{exo:m2:interest-rate-swaps:4}
The ten-year swap's DV01 is USD~82\,278 from the \omterm{def:m2:interest-rate-swaps:par}{annuity} and USD~82\,262 from bumping
and rebuilding the curve. Why do they differ, and which would you hedge with?
\end{exercise}

\begin{exercise}[$\star\star$]\label{exo:m2:interest-rate-swaps:5}
The 5y5y forward swap of \cref{fig:m2:interest-rate-swaps:buckets} has a par
rate of 4.295\%. Explain its bucketed DV01 and give the par swaps that hedge it.
\end{exercise}

\begin{exercise}[$\star\star$]\label{exo:m2:interest-rate-swaps:6}
Dealer A pays fixed to B, B pays fixed to C, C pays fixed to A, each USD~100
million of the same ten-year swap. What is the gross notional, the net risk of
each dealer, and the result of compression?
\end{exercise}

\begin{exercise}[$\star\star\star$]\label{exo:m2:interest-rate-swaps:7}
\emph{Coding.} With \texttt{firm\_curve}, bootstrap the curve of
\cref{ex:m2:interest-rate-swaps:ten} and report the four-year and eight-year
par rates and the ten-year discount factor.
\end{exercise}

\begin{exercise}[$\star\star\star$]\label{exo:m2:interest-rate-swaps:8}
\emph{Find the flaw.} ``The thirty-year \omterm{def:m2:interest-rate-swaps:spread}{swap spread} is negative, so the market
thinks banks are safer than the US government.'' Correct it.
\end{exercise}

\section{Problem: The Corporate Hedge}

\begin{problem}[{Weekend problem --- fixing a ten-year loan}]\label{pb:m2:interest-rate-swaps:1}
A company has borrowed USD~200 million for ten years at compounded SOFR plus
1.50\%. On 25 September 2026 it enters a ten-year \omterm{def:m2:interest-rate-swaps:payer}{payer swap} starting on 29
September at the par rate of the curve of \cref{ex:m2:interest-rate-swaps:ten}.

\medskip\noindent\textbf{Part I --- The swap.}
\begin{enumerate}
\item Give the fixed rate and the company's all-in fixed cost.
\item Give the swap's DV01 for the company, with its sign.
\item What is the swap worth to the company on the trade date?
\item Which pillar carries its risk?
\item Why does the swap start on 29 September and not on the trade date?
\end{enumerate}

\medskip\noindent\textbf{Part II --- Scenarios.}
\begin{enumerate}[resume]
\item All par rates rise by 50 basis points. What is the swap worth?
\item All fall by 50 basis points.
\item Why are the two not equal and opposite?
\item The curve steepens: $-20$ basis points at one year rising linearly to $+15$
at ten. Give the shocks at each pillar and the swap's value.
\item Why is the steepener worth almost exactly 15 times the DV01?
\end{enumerate}

\medskip\noindent\textbf{Part III --- The accounts.}
\begin{enumerate}[resume]
\item After the 50-basis-point fall the swap shows a loss. Has the company lost
money?
\item What does hedge accounting try to achieve, and why do treasurers care?
\item The bank quoted 4.07\% instead of the 4.05\% mid. What did the company pay
for the hedge, in dollars today?
\item What credit risk does each side take, and how is it managed?
\item What would a five-year swap have left unhedged?
\end{enumerate}

\medskip\noindent\textbf{Part IV --- Judgement.}
\begin{enumerate}[resume]
\item Why might a company prefer a swap to issuing a fixed-rate bond?
\item How does the bank hedge the swap it has just written?
\item Why is the company's hedge a \omterm{def:m2:interest-rate-swaps:payer}{payer swap} and a pension fund's a \omterm{def:m2:interest-rate-swaps:payer}{receiver
swap}?
\item State the \emph{named result}: the fixed rate, the DV01 of the hedge, and its
value after the steepening.
\item In one sentence: what does a swap exchange?
\end{enumerate}
\end{problem}

\section{Interview questions}

\begin{interviewq}[$\star$ \iqroles{trader, bank}]\label{iq:m2:interest-rate-swaps:1}
What is an \omterm{def:m2:interest-rate-swaps:swap}{interest-rate swap}, and why would a company enter one?
\end{interviewq}

\begin{interviewq}[$\star$ \iqroles{researcher, bank}]\label{iq:m2:interest-rate-swaps:2}
How do you price the \omterm{def:m2:interest-rate-swaps:swap}{floating leg} of an \omterm{def:m2:interest-rate-swaps:ois}{overnight index swap}?
\end{interviewq}

\begin{interviewq}[$\star\star$ \iqroles{researcher, developer}]\label{iq:m2:interest-rate-swaps:3}
Walk me through bootstrapping a discount curve from swap rates. What choices
do you make, and how do you test the result?
\end{interviewq}

\begin{interviewq}[$\star\star$ \iqroles{trader, researcher}]\label{iq:m2:interest-rate-swaps:4}
Why have thirty-year US \omterm{def:m2:interest-rate-swaps:spread}{swap spreads} been negative since 2008?
\end{interviewq}

\begin{interviewq}[$\star\star$ \iqroles{trader}]\label{iq:m2:interest-rate-swaps:5}
You are receiving fixed on a 5y5y forward swap. What is your risk, and how do
you hedge it with par swaps?
\end{interviewq}

\begin{interviewq}[$\star\star\star$ \iqroles{developer, researcher}]\label{iq:m2:interest-rate-swaps:6}
Your bucketed DV01s do not sum to the parallel DV01. List the possible causes.
\end{interviewq}
